\subsection{The exact defect at every level}
Keep the notation above and assume $t\ge1$. Put $d_0=q-1$ and $a=D/d_0=q^{t-1}$. All exponents occurring in $f_h(X)$ are $1\bmod d_0$, by induction in the recurrence for $f_h$. Consequently the reciprocal polynomial satisfies
\[
   X^{-\omega}\prod_{c\ne b} f_{h_c}(X)=F(X^{-d_0}),\qquad F\in\mathcal O[z],\quad F(0)=1 .
\]
In particular the loss $U_r-T_r$ is always divisible by $d_0$. Write $v_p$ for the ordinary $p$-adic valuation on integers, and set $e=v(p)$.

\paragraph{Theorem.}
There exists a bounded eventually periodic sequence $\eta_r$ of integers such that
\[
 T_r=\omega+d_0\{a(r-1)-\eta_r\},\qquad 0\le\eta_r\le (m-1)a.
\]
The following construction computes its exact values, including the initial ones, from a single finite ring independent of $r$.
Adjoin $\rho$ with $\rho^a=\pi$, put $N=ma$ and use
$S=\mathcal O[\rho]/(\rho^N)$. Here $v_\rho(\rho)=1$, with the convention that the valuation of a zero class in $S$ is $+\infty$. Set
\[
 G(z)=F(z/\rho),\qquad
 \frac1{G(z)}=\sum_{j\ge0}g_jz^j\quad\text{in }S[[z]],\qquad g_j=0\quad(j<0).
\]
Then $\eta_r$ is the least integer $i\in[0,(m-1)a]$ satisfying
\[
             v_\rho\bigl(g_{a(r-1)-i}\bigr)<a+i .
\]
\paragraph{Proof.}
The degree allowance established above gives $a v([z^j]F)\ge j$, so $G$ is integral. On reducing $G$ modulo $\rho$, only the constant terms and the first-order terms from the $m$ deepest factors survive. Thus $G\equiv(1+z^a)^m\bmod\rho$. If $H_j=[z^j]F^{-1}$, then $g_j$ is the reduction of $H_j\rho^{-j}$. At $j=a(r-1)-i\ge0$ the condition $v(H_j)<r$ is exactly $v_\rho(H_j\rho^{-j})<a+i$, since $v(H_j)$ is integral. All these cutoffs are at most $N$. The preceding theorem ensures there is a surviving $H_j$ in this range. The largest surviving $j$ gives the pole $\omega+d_0 j$, proving the formula.

Here is an explicit, possibly nonminimal, period. Put $L=mN$, $E=G-(1+z^a)^m$ in $S[z]$, and
\[
 C(z)=\sum_{k=0}^{N-1}(-E)^k(1+z^a)^{m(N-1-k)} .
\]
Since $E^N=0$, the identity $G^{-1}=C(z)/(1+z^a)^L$ holds in $S[[z]]$. For each class modulo $a$, and for $j\ge\deg C$, the coefficient $g_j$ is a finite linear combination over $S$ of shifted sequences $(-1)^u\binom{u+L-1}{L-1}$, indexed by $u=(j-k)/a$ for the relevant monomials $z^k$ in $C$. Let
\[
 b_0=\lceil m/e\rceil,\qquad
 A=2p^{\,b_0+\lfloor\log_p(\max\{1,L-1\})\rfloor}.
\]
Then $p^{b_0}=0$ in $S$ because $v_\rho(p)=ae$. For $1\le l\le L-1$, the identity $l\binom A l=A\binom{A-1}{l-1}$ gives $p^{b_0}\mid\binom A l$. Vandermonde's identity shows that the displayed signed binomial sequence has period dividing $A$ in $S$. This includes $L=1$, when evenness of $A$ suffices. Hence $g_{j+aA}=g_j$ for $j\ge\deg C$, and $\eta_{r+A}=\eta_r$ whenever $a(r-1)-(m-1)a\ge\deg C$.\hfill$\square$

\paragraph{A closed formula at depth one.}
If $t=1$, then $F(z)=(1+\pi z)^m$. Here the full answer simplifies to
\[
 T_r=\omega+(q-1)\max\left\{k\in\{0,\ldots,r-1\}:         k+e\,v_p\binom{k+m-1}{m-1}<r\right\}.
\]
Indeed the $k$th coefficient of $F^{-1}$ is $(-\pi)^k\binom{k+m-1}{m-1}$. For example, with $q=p=2$, $m=2$ and $e=1$, the maximum is $r-1$ for odd $r$ and $r-2$ for even $r$. One realizes this example at $b=0$ with directions $0,2,6$ over $\mathbb Z_2$.

\paragraph{Computation over the original local ring.}
For implementation the auxiliary root $\rho$ can be removed. Define $k_j=\lceil j/a\rceil$, $c_i=[z^i]F/\pi^{k_i}$ and $B_j=H_j/\pi^{k_j}$. Dividing the recurrence $[z^j](FF^{-1})=0$ by $\pi^{k_j}$ gives $B_0=1$ and, for $j>0$,
\[
 B_j=-\sum_{i=1}^{\min(j,\deg F)}c_i B_{j-i}\pi^{k_i+k_{j-i}-k_j}\pmod{\pi^m}.
\]
The exponents in this expression are zero or one. For $j=a(r-1)-i\ge0$, the survival test becomes $B_j\not\equiv0\bmod\pi^{\lfloor i/a\rfloor+1}$. This recurrence computes every level in the fixed ring $\mathcal O/\pi^m$.

At precision $r=2$ there is also a formula for arbitrary depths. Let $m_h=\#\{c\ne b:h_c=h\}$. The congruence for $f_h$ modulo $\pi^2$ gives
\[
 T_2=\omega+\max\bigl(\{0\}\cup
   \{q^h-q^{h-1}:h\ge1,\ p\nmid m_h\}\bigr).
\]
For $p=2$ and $q=4$, compare contact multisets $(0,2,2)$ and $(1,2,2)$. Both have $D=12$ and $m=2$, and neither reaches $U_2$. In the first case the loss is $12$. In the second the depth-one factor leaves a surviving pole increase of $3$, and the loss is $9$. Thus shallower contacts can determine the exact answer even at the first cancellation.

\paragraph{A classical special case.}
Over $K=\mathbb Q_p$ take $\pi=p$ and a single factor of depth $t\ge1$. In the multiplicative formal group its division polynomial is $(1+X)^{p^t}-1$. The condition $X^n/\Delta_b\in A_r[[X]]$ says exactly that $X^n=0$ in the quotient by this polynomial. Weierstrass division identifies the quotient with $(\mathbb Z/p^r)[T]/(T^{p^t}-1)$ via $X=T-1$. Our value $p^t+(r-1)(p-1)p^{t-1}$ agrees with the classical nilpotency index of the augmentation ideal of this cyclic group ring. See P. L. Clark and U. Schauz, \emph{Functional degrees and arithmetic applications, I}, Theorem~4.7 and its second proof, which recalls results of Weisman and Wilson (\url{https://arxiv.org/abs/2201.02763}). The multibranch problem involves products of division polynomials in the same variable, including repeated factors.
